\begin{lemma}
Fix \(n,m,p\), an event \(E\) in the two-bite sample space, and a
nonnegative constant \(C\).  Assume \(0\le p\le1\).  Suppose that for every
blue base graph \(G_B\) and every blue row map
\(\rho:\operatorname{Fin}n\to\operatorname{Fin}m\), the conditional
probability of \(E\) given the cell
\[
\omega.2.1=G_B,\qquad v\mapsto (\omega.2.2(v))_2=\rho(v)
\]
is at most \(C\).  Then the unconditional two-bite probability of \(E\) is
at most \(C\).
\end{lemma}
\begin{proof}
For a sample \(\omega\), let
\[
\Pi_B(\omega)=\left(\omega.2.1,\ v\mapsto(\omega.2.2(v))_2\right).
\]
The values of \(\Pi_B\), namely pairs \((G_B,\rho)\), index a finite
partition of the two-bite sample space: every sample has exactly one blue
base graph and exactly one blue-coordinate row map.  Let
\[
C_{G_B,\rho}=\{\omega:\Pi_B(\omega)=(G_B,\rho)\}
\]
be the corresponding cell, and write \(\mu(A)\) for
\(\operatorname{TwoBiteEventProbability}(A)\).  Expanding the finite weighted
sum by fibers of \(\Pi_B\) gives
\[
\mu(E)=\sum_{G_B,\rho}\mu(E\cap C_{G_B,\rho}).
\]

By \(0\le p\le1\), \(\noderef{TwoBiteSampleWeightNonnegative}\) says that
every two-bite sample weight is nonnegative.  Fix a cell
\(C_{G_B,\rho}\).  If \(\mu(C_{G_B,\rho})=0\), then
\(\mu(E\cap C_{G_B,\rho})=0\) as well: the latter event is a subset of the
cell, and a finite sum of nonnegative summands over a subset of a zero-mass
set is again zero.  This is also the zero-mass convention in
\(\noderef{TwoBiteConditionalProbability}\), where the conditional
probability on such a cell is defined to be \(0\).

If \(\mu(C_{G_B,\rho})\ne0\), nonnegativity gives
\(\mu(C_{G_B,\rho})>0\), and the conditional probability on the cell is
\[
\frac{\mu(E\cap C_{G_B,\rho})}{\mu(C_{G_B,\rho})}.
\]
The hypothesis bounds this ratio by \(C\).  Multiplying by the positive cell
mass gives
\[
\mu(E\cap C_{G_B,\rho})\le C\,\mu(C_{G_B,\rho}).
\]
The same inequality also holds in the zero-mass case because both sides are
zero or nonnegative and \(C\ge0\).

Summing over all finite blue cells gives
\[
\mu(E)
\le C\sum_{G_B,\rho}\mu(C_{G_B,\rho}).
\]
The cells partition the sample space, so the sum of their masses is the total
two-bite mass.  By \(\noderef{TwoBiteEventProbabilityTotalMassBound}\), this
total mass is at most \(1\).  Since \(C\ge0\), we obtain
\[
\mu(E)\le C\cdot 1=C.
\]
This is the desired unconditional bound.
\end{proof}
